\section{Navigation: MPPI, the costmap, and how runs are measured}
\label{sec:nav}

\subsection{The MPPI controller}

Nav2 follows its planned path with MPPI (model predictive path integral). Every
control cycle it samples $K$ sequences of commands $(v_x, v_y, \omega)$ over a
horizon of $N$ steps of $\Delta t$, rolls each through a kinematic model of the
boat (Omni: the X thrusters can sway), scores each rollout with a sum of critics
$S_k$ (distance from the path, heading along it, nearness to obstacles, progress to
the goal, turning), and blends them with weights
\[
  w_k = \frac{e^{-(S_k - S_{\min})/\lambda}}{\sum_j e^{-(S_j - S_{\min})/\lambda}},
  \qquad \bm u = \sum_k w_k\,\bm u_k .
\]
The temperature $\lambda$ sets how much it averages: a high $\lambda$ blends many
samples, a low one commits to the best few. A floating boat's odometry is noisy
(waves) and its response lags, and averaging over many samples then washes out
the corrections; $\lambda = 0.1$ commits.

\paragraph{Timing.} After each cycle MPPI shifts its best sequence one step,
assuming a cycle lasts $\Delta t$. So $\Delta t$ must equal the control period:
$\Delta t = 1/f$. At $f = 20$\,Hz on a loaded computer the loop slipped to 11--17\,Hz,
the shifted plan fell out of step with the boat, and it drove in circles; at
$f = 10$\,Hz with $\Delta t = 0.1$\,s and $N=28$ (the same 2.8\,s horizon) it holds
its rate and drives straight.

\paragraph{Path horizon.} MPPI scores rollouts against the part of the path
within the prune distance $D$. A rollout runs up to $v_{\max} N\Delta t$; if that
exceeds $D$ the rollouts run off the end of the path, where only weak critics
apply, and near the end of a path the boat circles instead of arriving. Hence
$D \ge 1.2\, v_{\max} N \Delta t$.

\subsection{The costmap over open water}

A LiDAR over water returns almost nothing: the surface does not reflect the beam
back and there is nothing behind it. An obstacle grid clears a cell only when a
ray passes through it to a return, so a mark made once (a buoy, a gust-tilted
ring of points) would never be cleared. The obstacle layer therefore uses voxels
that decay: each mark lives a few seconds ($3$\,s in the local map, $15$\,s in the
global one, which must remember a buoy behind the boat) unless seen again. And a
VLP-16's rings are $2^\circ$ apart: a buoy $0.15$\,m above the water, $2$\,m away,
fits between them. The camera keeps such buoys in view, so the tracker's objects
are marked as small rings of points as a second source.

\paragraph{Inflation and the footprint.} The hull is a $1.06\times0.90$\,m
rectangle, padded by $0.1$\,m for tracking error at speed and for the scatter in
where a buoy is seen. Its padded corners lie
$r_c = \sqrt{0.63^2 + 0.55^2} \approx 0.84$\,m from the centre. MPPI scores
most trajectory points by the cost under the centre alone, and checks the whole
rectangle only where that cost says a corner could reach something, which is
the cost found at $r_c$. If the inflation ends short of $r_c$ that test can never
fire, and a corner can swing across a buoy the centre keeps clear of.

The inflation reaches further, to $1.5$\,m, and falls off gently, as
$252\,e^{-3(d - r_i)}$ beyond the inscribed radius $r_i = 0.55$\,m. Had it ended
at half a gate, the middle of any gap between two buoys a little wider than a
gate would cost nothing, and the planner, going from one task to the next, would
take the shortest line through the buoy field, sliding past a buoy with only the
footprint check's few centimetres to spare. Where a buoy is seen moves by
$0.1$ to $0.3$\,m while the boat turns, which is more than that. With the wider
inflation the gap costs something, and open water around the field is cheaper.
A gate's middle still has to be passable: in a 6\,ft ($1.83$\,m) gate it is
$0.91$\,m from each buoy and costs $252\,e^{-3(0.91-0.55)} \approx 84$, well
short of the inscribed cost, so the boat goes through when its goal lies beyond.

\subsection{Measuring a run}

\paragraph{Cross-track error.} The route is a polyline through the waypoints,
starting where the boat was when the goal was sent. For a segment from $\bm a$
to $\bm b$ and a position $\bm p$, project onto the segment and clamp to its ends:
\[
  \lambda = \operatorname{clip}\!\left(\frac{(\bm p-\bm a)\cdot(\bm b-\bm a)}{\|\bm b-\bm a\|^2},\,0,\,1\right),\qquad
  d = \left\|\bm p - (\bm a + \lambda(\bm b-\bm a))\right\|.
\]
The cross-track error of $\bm p$ is the smallest $d$ over all segments. Clamping
matters at corners: without it the distance would be to the infinite line, which
can be small far beyond the end of a leg. The procedure reports the mean, RMS and
maximum, the closest approach to each waypoint, and the travelled path length
against the route length.

\paragraph{Station keeping.} With samples $\bm p_k$ at times $t_k$ around a
centre $\bm c$, the radii $\rho_k = \|\bm p_k - \bm c\|$ give the RMS and maximum
radius, and the fraction of time inside the tolerance (each sample weighted by its
interval). A slow steady push that the controller does not cancel can stay inside
the tolerance for a long time, so the drift is measured directly as the
least-squares slope of each coordinate:
\[
  (\alpha_x, v_x) = \argmin_{\alpha,v} \sum_k \left(x_k - \alpha - v\,(t_k - t_0)\right)^2,
\]
and likewise for $y$. The drift speed is $\sqrt{v_x^2+v_y^2}$ and its heading
$\operatorname{atan2}(v_y, v_x)$.
\file{mhseals\_testing/analysis/nav\_metrics.py}
